% cfr_connection.tex - Connection to CFR
\documentclass[11pt]{article}
\usepackage{amsmath, amssymb, amsthm}

\title{No-Regret Learning in CFR and Market Making}
\author{}
\date{}

\begin{document}
\maketitle

\section{CFR}
Counterfactual Regret Minimization maintains a regret minimizer at each
information set of an extensive-form game. If each local regret is sublinear,
then total regret is also sublinear:
\[
R_T \le \sum_{I} R_T^I.
\]
The average strategy therefore approaches a Nash equilibrium in two-player
zero-sum games.

\section{Spread Selection}
In the market-making experiment, the repeated decision is the spread choice.
The Exp3 learner receives a reward equal to spread revenue plus inventory PnL
minus risk penalties. Its regret is measured against the best fixed spread in
hindsight:
\[
R_T = \max_s \sum_{t=1}^T r_t(s)-\sum_{t=1}^T r_t(s_t).
\]

\section{Decomposition}
The simulator also records economic components
\[
\text{PnL}_t
=
\text{spread}_t
- \text{adverse\_selection}_t
- \text{inventory\_loss}_t.
\]
This is not a separate learning guarantee; it is an accounting identity used to
inspect where rewards come from. The shared structure with CFR is the use of
local regret minimization to control a global performance criterion.

\end{document}
